\documentclass[12pt, a4paper]{ctexart}
\usepackage[margin=2cm]{geometry}
\usepackage{libertine}

\usepackage{titlesec}
\usepackage{zhnumber}
\titleformat*{\section}{\Large\bfseries\raggedright}
\renewcommand{\thesection}{\zhnum{section}}
\renewcommand{\thesubsection}{\arabic{subsection}}

\setcounter{secnumdepth}{4}
\renewcommand{\theparagraph}{\thesubsubsection.\arabic{paragraph}}
\renewcommand{\paragraph}[1]{
    \refstepcounter{paragraph}
    {\bfseries\theparagraph\quad#1}\par
    \vspace{2pt}
    \noindent
}

\usepackage{enumitem}
\setlist[enumerate]{itemsep=2pt, parsep=0pt, partopsep=0pt, topsep=2pt}
\setlist[itemize]{itemsep=2pt, parsep=0pt, partopsep=0pt, topsep=2pt}
\linespread{1.2}

\usepackage{graphicx}

\usepackage{minted}
\setminted{
    breaklines=true, escapeinside=||, fontsize=\small, frame=lines,
    mathescape=$$, numbers=none, style=bw, tabsize=4
}

\usepackage{siunitx}

\begin{document}
    
\pagestyle{plain}
\thispagestyle{empty}

\noindent
\begin{tabular*}{\textwidth}{l @{\extracolsep{\fill}} r @{\extracolsep{6pt}} l}
    \LARGE{\textbf{实验报告}} & 计算机网络 & \textit{Computer Networking} \\
\end{tabular*}\\\\
\begin{tabular*}{\textwidth}{l l}
    \textbf{报告标题: } & \textbf{VLAN 及 VLAN 间路由选择} \\
    \textbf{学号: } & 19240212 \\
    \textbf{姓名: } & 华博文 \\
    \textbf{日期: } & 2026 年 1 月 5 日 \\
\end{tabular*}\\
\rule[2ex]{\textwidth}{2pt}

\section{实验目的}

本实验旨在掌握基于 IPMininet 的 VLAN（虚拟局域网）配置方法，理解 VLAN 实现二层网络隔离、同 VLAN 主机互通的核心原理；通过 IPMininet 搭建 ``1 台交换机 + 3 台主机'' 的网络拓扑，完成主机 IP 子网掩码配置、Open vSwitch（OVS）交换机安全模式调整、VLAN 端口标签划分，以及验证不同 VLAN 下主机的连通性，同时熟悉 IPMininet 仿真工具与 OVS 命令在虚拟网络实验中的操作流程。

\section{实验内容简要描述}

\begin{enumerate}
    \item \textbf{IPMininet 仿真网络启动}：在 Linux 环境启动 Open vSwitch 服务，通过 \mintinline{bash}`mn` 命令启动包含 1 台交换机、3 台主机的单拓扑仿真网络；
    
    \item \textbf{主机与交换机配置}：打开各主机终端，调整网卡的子网掩码至实验规划的 \mintinline{text}`255.255.255.0`，匹配网段要求；打开交换机终端，查看 OVS 状态并退出安全模式，确保交换机具备独立转发能力；
    
    \item \textbf{VLAN 配置}：通过 OVS 命令查看主机与交换机端口的 MAC 映射关系，为交换机端口配置 VLAN 标签（划分 VLAN 100、200）；
    
    \item \textbf{主机连通性验证}：测试同 VLAN、不同 VLAN 主机的 \mintinline{bash}`ping` 连通性，验证 VLAN 的二层隔离效果。
\end{enumerate}

\section{实验步骤与结果分析}

\subsection{IPMininet 仿真网络启动}

右键打开 Linux 终端，执行命令启动 OVS 服务（为交换机提供交换功能支持）：

\begin{minted}{bash}
ovs-vswitchd --pidfile --detach
\end{minted}

在终端中执行命令启动包含 1 台交换机（s1）、3 台主机（h1、h2、h3）的仿真网络：

\begin{minted}{bash}
sudo mn --topo=single,3 --controller=remote
\end{minted}

终端依次显示 ``Creating network'' ``Adding hosts'' ``Adding switches'' 等日志，最终进入 Mininet CLI 界面，确认网络启动完成。

\begin{figure}[H]
    \centering
    \includegraphics[width=\textwidth]{./img/screenshot_0000.png}
\end{figure}

\subsection{主机与交换机配置}

在 Mininet CLI 中打开 h1 的终端：

\begin{minted}{bash}
xterm h1
\end{minted}

执行 \mintinline{bash}`sudo ifconfig` 查看默认 IP 信息，发现 eth0 网卡的子网掩码为 \mintinline{text}`255.0.0.0`（不符合实验规划），执行命令修改子网掩码为 \mintinline{text}`255.255.255.0`：

\begin{minted}{bash}
sudo ifconfig h1-eth0 netmask 255.255.255.0
\end{minted}

执行 \mintinline{bash}`sudo ifconfig` 确认配置生效（h1-eth0 的 netmask 显示为 \mintinline{text}`255.255.255.0`）：

\begin{figure}[H]
    \centering
    \includegraphics[width=\textwidth]{./img/screenshot_0001.png}
\end{figure}

同理，在 Mininet CLI 中打开 h2、h3 的终端，执行相同命令修改子网掩码：

\begin{figure}[H]
    \centering
    \includegraphics[width=\textwidth]{./img/screenshot_0002.png}
\end{figure}

在 Mininet CLI 中打开交换机 s1 的终端：

\begin{minted}{bash}
xterm s1
\end{minted}

执行 \mintinline{bash}`ovs-vsctl show` 查看交换机状态，发现 fail\_mode 为 ``secure''（表示交换机依赖 SDN 控制器转发，需退出该模式）。

执行命令退出安全模式：

\begin{minted}{bash}
ovs-vsctl del-fail-mode s1
\end{minted}

再次执行 \mintinline{bash}`ovs-vsctl show`，确认 fail\_mode 已移除，交换机切换为普通转发模式：

\begin{figure}[H]
    \centering
    \includegraphics[width=\textwidth]{./img/screenshot_0003.png}
\end{figure}

\subsection{VLAN 配置}

在 s1 的终端中，执行命令查看交换机端口与主机 MAC 的映射关系：

\begin{minted}{bash}
ovs-appctl fdb/show s1
\end{minted}

结合各主机终端中 \mintinline{bash}`sudo ifconfig` 显示的 MAC 地址，确认端口对应关系：s1-eth1 连接 h1、s1-eth2 连接 h2、s1-eth3 连接 h3。

\begin{figure}[H]
    \centering
    \includegraphics[width=\textwidth]{./img/screenshot_0004.png}
\end{figure}

执行命令为交换机端口配置 VLAN 标签（h1、h2 划入 VLAN 100，h3 划入 VLAN 200）：

\begin{minted}{bash}
ovs-vsctl set port s1-eth1 tag=100
\end{minted}

\begin{minted}{bash}
ovs-vsctl set port s1-eth2 tag=100
\end{minted}

\begin{minted}{bash}
ovs-vsctl set port s1-eth3 tag=200
\end{minted}

执行 \mintinline{bash}`ovs-vsctl show` 确认端口 VLAN 标签配置生效（端口后显示 ``tag: 100'' 或 ``tag: 200''）：

\begin{figure}[H]
    \centering
    \includegraphics[width=\textwidth]{./img/screenshot_0005.png}
\end{figure}

\subsection{主机连通性验证}

在 Mininet CLI 中打开 h1 的终端，执行 \mintinline{bash}`ping 10.0.0.2` 命令 ping 同 VLAN 的 h2：终端显示 ``64 bytes from 10.0.0.2'' 的回复，丢包率为 0\%，说明同 VLAN 主机互通。

继续在 h1 终端中 ping 不同 VLAN 的 h3：终端显示 ``Destination Host Unreachable''，丢包率为 100\%，说明不同 VLAN 主机隔离。

同理，在 h2 终端中 ping h3，结果同样显示无法连通，进一步验证了 VLAN 的二层隔离作用。

\begin{figure}[H]
    \centering
    \includegraphics[width=\textwidth]{./img/screenshot_0006.png}
\end{figure}

\section{实验中遇到的问题及体会}

通过本次实验，我直观理解了 VLAN 的核心作用：在二层网络中划分逻辑网段，实现同 VLAN 互通、不同 VLAN 隔离，避免广播域过大的问题。同时，OVS 作为虚拟交换机的配置方式（如端口标签、模式调整），让我熟悉了软件定义交换的基本操作。

实验也深化了 ``二层隔离'' 的认知：VLAN 是在数据链路层对流量进行划分，即使主机处于同一 IP 网段，不同 VLAN 的主机也无法直接通信，这与三层路由的网段互通逻辑形成了区别。这种分层隔离的设计，是网络安全与性能优化的常用手段。

最后，IPMininet 的虚拟环境提供了低成本的 VLAN 实践载体，可快速调整拓扑与配置、反复验证效果。这一过程不仅提升了我的 Linux 命令与 OVS 配置熟练度，更培养了从端口标签、连通性结果倒推 VLAN 配置错误的排查思维，为后续 VLAN 间路由等高级技术的学习奠定了基础。

\end{document}